\section{Introduction}

Options market makers operate under regulatory mandates to maintain delta-neutral books despite accumulating gamma exposure~\citep{sec15c31}. When aggregate gamma exposure turns substantially negative, dealers face forced pro-cyclical hedging—selling strength and purchasing weakness—thereby amplifying price movements through coordinated rebalancing flows~\citep{ni2005stock,garleanu2009demand}. Practitioners commonly summarize these dynamics through net gamma exposure (GEX) metrics, compressing rich strike-level options data into scalar values for tractability. Yet a critical question remains: does this parametric compression discard structural signal relevant for detecting dealer constraints?

We demonstrate that it does. Large language models processing raw strike-level data—open interest, implied volatility, and bid-ask spreads per strike—achieve 92.3\% detection of dealer hedging patterns, outperforming identical models given pre-calculated Net GEX by 30.8 percentage points. This finding establishes that standard industry metrics represent lossy compression of actionable structural information. The options surface contains signal about dealer positioning that scalar GEX values cannot capture: asymmetrical strike concentrations, volatility skew patterns, and spatial distribution of open interest relative to spot price.

This result emerges from a rigorous validation methodology we term \textit{obfuscation testing}. Presenting a model with ``SPY on January 27, 2021'' alongside substantial negative gamma risks measuring memorization of the GameStop squeeze rather than comprehension of underlying dealer mechanics. This training data contamination problem pervades machine learning evaluation in finance, where temporal context enables recall rather than reasoning. Obfuscation testing eliminates all memorizable information while preserving structural relationships: converting dates to generic labels (``Day T+0''), anonymizing ticker symbols (``SPY'' → ``INDEX\_1''), and removing economic event references. Successful pattern detection under these conditions validates genuine structural analysis rather than statistical correlation with training data.

The proliferation of zero-days-to-expiration (0DTE) options establishes a particularly rigorous validation environment. Trading volume in 0DTE contracts expanded from negligible levels pre-2022 to exceeding 50\% of total SPX options volume by 2024~\citep{cboe2023}, generating a persistent negative gamma regime where dealers maintain net short gamma across 95.6\% of trading days in our sample. This structural transition from historical alternating regimes~\citep{ni2005stock,garleanu2009demand} to sustained single-regime dynamics presents a demanding test case: unlike conventional pattern recognition exploiting regime variation, our methodology must identify the underlying constraint mechanism within a uniform environment.

Our methodology advances financial data science through four contributions. First, we establish that unstructured feature extraction from raw options chains recovers signal lost in parametric GEX compression—a 30.8 percentage point detection improvement with direct implications for risk system design. Second, we introduce obfuscation testing as a validation framework for distinguishing genuine structural pattern detection from training data memorization. Third, we quantify prompt sensitivity effects: including regime labels (``NEGATIVE\_GAMMA'') inflates detection from 71.5\% to 100\%, revealing how contextual hints generate misleading performance metrics. Fourth, we establish that detected patterns function as risk management signals rather than alpha generators: stable detection capability (68-74\% quarterly) persists as economic profitability collapses (Sharpe 1.8 → 0.1), and detection days exhibit 33\% lower volatility than baseline—proving identification of suppression mechanisms rather than high-volatility events.

Testing this framework on 242 trading days throughout 2024, we evaluate three manifestations of dealer hedging constraints: gamma positioning, stock pinning, and 0DTE hedging. Using fully unbiased prompts with obfuscated data, we achieve 71.5\% average detection rate, with 91.2\% of detected patterns materializing in forward returns. The WHO→WHOM→WHAT causal framework requires articulation of economic actors, affected parties, and forced mechanisms, ensuring detected patterns reflect genuine market microstructure dynamics.

These findings have implications for both risk management and AI deployment in finance. For practitioners, our results demonstrate that standard GEX metrics underutilize available structural information, suggesting opportunities to improve risk system inputs through unstructured feature extraction. For researchers, obfuscation testing offers a rigorous framework for validating pattern detection in domains where training data contamination is a concern.

The remainder of this paper is organized as follows. Section II reviews related work. Section III presents our methodology. Section IV describes the experimental setup. Section V reports results. Section VI discusses implications and limitations. Section VII concludes.
